% ============================================================
\chapter{Delta Hedging and Risk Management}
\label{ch:hedging}
% ============================================================

The Black-Scholes model does not merely give a price: it tells you how to \emph{hedge}. The central idea is that the risk of an option can be eliminated by continuously holding the right quantity of the underlying asset---the \emph{delta}. By adjusting this position at every instant, the portfolio becomes insensitive to market fluctuations. But perfect hedging is an ideal: in practice, one can only rebalance at discrete intervals, transaction costs accumulate, and volatility is never constant. Understanding these limitations is as important as understanding the theory.

\section{Introduction}

Hedging is the activity of reducing or eliminating risk associated with a
derivatives position. This chapter develops the theory and practice of
delta hedging, gamma hedging, and analyses the costs and limitations of
dynamic hedging.

\section{Delta Hedging}

\subsection{Principle}

\begin{definition}[Delta Hedging]
\textbf{Delta hedging} consists of maintaining a portfolio
$\Pi_t = V_t - \Delta_t S_t$
locally insensitive to changes in $S$, where
$\Delta_t = \frac{\partial V}{\partial S}(S_t, t)$.
\end{definition}

\begin{theorem}[Self-Financing Delta-Hedged Portfolio]
Under the Black--Scholes framework with continuous trading, the delta-hedged
portfolio is riskless and perfectly replicates the derivative. Its value
satisfies:
\[
    d\Pi_t = r\Pi_t\,dt.
\]
\end{theorem}

\begin{intuition}[title=\textbf{Discrete-Time Hedging}]
In practice, rebalancing is discrete (every $\Delta t$ time units). The
hedging error arises because $\Delta$ changes between rebalancings. Higher
frequency improves the hedge but increases transaction costs.
\end{intuition}

\subsection{Discrete-Time Hedging Error}

\begin{proposition}[Hedging Error]
For a European call hedged with delta rebalancing at dates
$t_0, t_1, \ldots, t_N$, the hedging error is:
\[
    \epsilon = \sum_{i=0}^{N-1} \frac{1}{2}\Gamma(S_{t_i}, t_i)
    S_{t_i}^2\bigl[(\Delta S_{t_i})^2 - \sigma^2 S_{t_i}^2 \Delta t\bigr]
    \cdot e^{r(T-t_{i+1})},
\]
where $\Delta S_{t_i} = S_{t_{i+1}} - S_{t_i}$.
\end{proposition}

The hedging error has zero expectation under $\mathbb{Q}$ but non-zero
variance, proportional to $\Delta t$.

\begin{codeblock}[title=\textbf{Discrete Delta-Hedging Simulation}]
\begin{minted}{python}
import numpy as np
from scipy.stats import norm

def bs_delta(S, K, T, r, sigma):
    """Delta of a European call."""
    if T < 1e-10:
        return 1.0 if S > K else 0.0
    d1 = (np.log(S/K) + (r + 0.5*sigma**2)*T) / (sigma*np.sqrt(T))
    return norm.cdf(d1)

def bs_price(S, K, T, r, sigma):
    """Black-Scholes call price."""
    if T < 1e-10:
        return max(S - K, 0)
    d1 = (np.log(S/K) + (r + 0.5*sigma**2)*T) / (sigma*np.sqrt(T))
    d2 = d1 - sigma*np.sqrt(T)
    return S*norm.cdf(d1) - K*np.exp(-r*T)*norm.cdf(d2)

def simulate_delta_hedge(S0, K, T, r, sigma, N_rebal, M=10000):
    """Simulate delta hedging and compute hedging errors."""
    dt = T / N_rebal
    errors = np.zeros(M)

    for m in range(M):
        S = S0
        V0 = bs_price(S, K, T, r, sigma)
        delta = bs_delta(S, K, T, r, sigma)
        cash = V0 - delta * S

        for i in range(N_rebal):
            tau = T - (i+1)*dt
            Z = np.random.standard_normal()
            S_new = S * np.exp((r - 0.5*sigma**2)*dt
                               + sigma*np.sqrt(dt)*Z)
            cash *= np.exp(r * dt)
            if tau > 1e-10:
                delta_new = bs_delta(S_new, K, tau, r, sigma)
            else:
                delta_new = 1.0 if S_new > K else 0.0
            cash -= (delta_new - delta) * S_new
            delta = delta_new
            S = S_new

        portfolio = cash + delta * S
        payoff = max(S - K, 0)
        errors[m] = portfolio - payoff

    return errors

np.random.seed(42)
for N in [10, 50, 252, 1000]:
    errs = simulate_delta_hedge(100, 100, 1.0, 0.05, 0.20, N)
    print(f"N={N:4d}: mean={np.mean(errs):.4f}, "
          f"std={np.std(errs):.4f}")
\end{minted}
\end{codeblock}

\section{Gamma Hedging}

\begin{definition}[Gamma Hedging]
\textbf{Gamma hedging} makes the portfolio insensitive not only to small
changes in $S$ (delta neutral) but also to larger changes (gamma neutral).
A second instrument (typically another option) is used to neutralise gamma.
\end{definition}

\begin{proposition}[Constructing a Gamma-Neutral Portfolio]
Let $V$ be the derivative to hedge and $O$ a hedging instrument with gamma
$\Gamma_O$. The position in $O$ needed to neutralise gamma is:
\[
    w_O = -\frac{\Gamma_V}{\Gamma_O}.
\]
After this hedge, the position in the underlying is adjusted to restore
delta neutrality:
\[
    \Delta_{\text{total}} = \Delta_V + w_O \Delta_O.
\]
\end{proposition}

\begin{example}
A trader is short a call ($\Delta = 0.55$, $\Gamma = 0.03$). They have
access to a call with different maturity ($\Delta_O = 0.40$,
$\Gamma_O = 0.05$).

Gamma-neutral position: $w_O = -(-0.03)/0.05 = 0.60$.

New delta: $-0.55 + 0.60 \times 0.40 = -0.31$.

Buy $0.31$ units of the underlying to become delta neutral.
\end{example}

\section{Vega Hedging}

\begin{definition}[Vega Hedging]
\textbf{Vega hedging} aims to neutralise the portfolio's sensitivity to
volatility. Since the underlying's vega is zero, options must be used.
\end{definition}

\begin{remark}
In practice, vega hedging is more important than gamma hedging because
implied volatility movements can generate substantial profit and loss (P\&L).
\end{remark}

\section{P\&L Decomposition}

\begin{theorem}[P\&L Decomposition of a Delta-Hedged Portfolio]
The instantaneous P\&L of a delta-hedged portfolio is:
\[
    d\text{P\&L} = \frac{1}{2}\Gamma S^2\bigl[(dS/S)^2 - \sigma^2\,dt\bigr]
                 + \mathcal{V}\,d\sigma_{\text{impl}} + \cdots
\]
The first term is the \textbf{gamma P\&L} (profitable when realised
volatility exceeds implied volatility), the second is the \textbf{vega P\&L}.
\end{theorem}

\begin{keyformulas}[title=\textbf{Taylor Approximation of P\&L}]
For a general portfolio, the Taylor approximation gives:
\[
    \Delta\text{P\&L} \approx \Delta\cdot\delta S
    + \frac{1}{2}\Gamma(\delta S)^2 + \Theta\,\delta t
    + \mathcal{V}\,\delta\sigma + \rho\,\delta r.
\]
\end{keyformulas}

\section{Transaction Costs and Optimal Hedging}

\begin{theorem}[Optimal Rebalancing Frequency (Leland)]
In the presence of proportional transaction costs $k\abs{\delta S}$,
Leland (1985) shows that optimal hedging uses a modified volatility:
\[
    \hat{\sigma}^2 = \sigma^2\left(1 + \sqrt{\frac{2}{\pi}}
    \frac{k}{\sigma\sqrt{\Delta t}}\right).
\]
\end{theorem}

\begin{proposition}[Whalley--Wilmott Model]
An alternative approach rebalances only when delta deviates sufficiently
from its target. The optimal no-trade band is:
\[
    \Delta_{\text{target}} \pm \left(\frac{3k\,e^{-r\tau}\Gamma^2 S^2}
    {2\lambda}\right)^{1/3},
\]
where $\lambda$ represents the trader's risk aversion.
\end{proposition}

\begin{codeblock}[title=\textbf{Hedging with Transaction Costs}]
\begin{minted}{python}
def hedge_with_costs(S0, K, T, r, sigma, N_rebal, k=0.001, M=5000):
    """Delta hedging with proportional transaction costs."""
    dt = T / N_rebal
    total_costs = np.zeros(M)
    hedge_errors = np.zeros(M)

    for m in range(M):
        S = S0
        V0 = bs_price(S, K, T, r, sigma)
        delta = bs_delta(S, K, T, r, sigma)
        cash = V0 - delta * S
        cost = k * abs(delta) * S

        for i in range(N_rebal):
            tau = T - (i+1)*dt
            Z = np.random.standard_normal()
            S_new = S * np.exp((r-0.5*sigma**2)*dt
                               + sigma*np.sqrt(dt)*Z)
            cash *= np.exp(r*dt)

            delta_new = bs_delta(S_new, K, max(tau, 1e-10), r, sigma)
            trade = delta_new - delta
            cost += k * abs(trade) * S_new
            cash -= trade * S_new
            delta = delta_new
            S = S_new

        portfolio = cash + delta * S
        payoff = max(S - K, 0)
        hedge_errors[m] = portfolio - payoff
        total_costs[m] = cost

    print(f"Average cost: {np.mean(total_costs):.4f}")
    print(f"Error std:    {np.std(hedge_errors):.4f}")
    return hedge_errors, total_costs

np.random.seed(42)
hedge_with_costs(100, 100, 1.0, 0.05, 0.20, 252, k=0.001)
\end{minted}
\end{codeblock}

\section{Hedging in Incomplete Markets}

\begin{remark}
When the market is incomplete (stochastic volatility, jumps), perfect
hedging is impossible. Alternative approaches include:
\begin{itemize}
    \item \textbf{Minimum variance hedging}: minimise $\Var(\Pi_T - H)$
          over admissible strategies.
    \item \textbf{Super-replication}: find the minimal cost of a portfolio
          that dominates the payoff.
    \item \textbf{Quantile hedging}: control the probability of loss.
\end{itemize}
\end{remark}

\section{Exercises}

\begin{exercise}[Delta-Hedging Simulation]
\begin{enumerate}
    \item Simulate delta hedging of a European call for
          $N = 10, 50, 252, 1000$ rebalancings.
    \item Plot the histogram of hedging errors for each $N$.
    \item Verify that the error variance decreases as $O(1/N)$.
\end{enumerate}
\end{exercise}

\begin{exercise}[Gamma Hedging]
\begin{enumerate}
    \item Construct a delta-gamma neutral portfolio from a short ATM call,
          a long OTM call, and the underlying.
    \item Simulate the P\&L over one week and compare with the delta-only
          hedged portfolio.
\end{enumerate}
\end{exercise}

\begin{exercise}[P\&L Analysis]
\begin{enumerate}
    \item Decompose the daily P\&L of a delta-hedged portfolio into delta,
          gamma, theta, and vega components.
    \item Show that for a delta-neutral portfolio,
          $\text{P\&L} \approx \frac{1}{2}\Gamma S^2(\sigma_r^2 - \sigma_i^2)\,dt$
          where $\sigma_r$ is realised volatility and $\sigma_i$ is implied
          volatility.
\end{enumerate}
\end{exercise}

\begin{exercise}[Transaction Costs]
\begin{enumerate}
    \item Implement Leland hedging with modified volatility.
    \item Compare total cost and hedging error with the standard method
          for different values of $k$.
    \item Implement the Whalley--Wilmott no-trade band and analyse the
          cost/risk trade-off.
\end{enumerate}
\end{exercise}
